% Part: lambda-calculus
% Chapter: syntax
% Section: free-variables

\documentclass[../../../include/open-logic-section]{subfiles}

\begin{document}

\olfileid{lam}{syn}{fv}
\olsection{المتغيرات الحرة}

يتناول حساب لامبدا الدوال، وتجريد لامبدا هو الكيفية التي تنشأ بها الدوال.
فالحدس أن~$\lambd[x][M]$ هو الدالة التي تعطي قيمها~$M$ عندما يسند وسيط
الدالة إلى~$x$. لكن ليست كل واقعة لـ$x$ في~$M$ ذات صلة: فإذا احتوى~$M$
تجريدًا آخر~$\lambd[x][N]$، كانت وقائع~$x$ في~$N$ ذات صلة
بـ$\lambd[x][N]$ لا بـ$\lambd[x][M]$. ولذلك فإن تجريد لامبدا~$\lambd[x]$
داخل~$\lambd[x][M]$ \emph{يربط} وقائع~$x$ في~$M$ التي لم يربطها تجريد
لامبدا آخر بالفعل، أي الوقائع \emph{الحرة} لـ$x$ في~$M$.

\begin{defn}[المجال]
إذا وقع~$\lambd[x][M]$ داخل حد~$N$، فإن الواقعة المقابلة لـ$M$ هي
\emph{مجال}~$\lambd[x]$.
\end{defn}


\begin{defn}[الواقعة الحرة والمقيدة]
  تكون واقعة المتغير~$x$ في حد~$M$ \emph{حرة} إذا لم تكن في مجال
  $\lambd[x]$، وتكون \emph{مقيدة} في غير ذلك. وتكون واقعة المتغير~$x$ في
  $\lambd[x][M]$ مقيدة بـ$\lambd[x]$ الابتدائية إذا وفقط إذا كانت واقعة
  $x$ في~$M$ حرة.
\end{defn}

\begin{ex}
  في~$\lambd[x][x y]$ يقع كل من~$x$ و$y$ في مجال~$\lambd[x]$، ولذلك
  يكون~$x$ مقيدًا بـ$\lambd[x]$. ولما كان~$y$ لا يقع في مجال أي
  $\lambd[y]$، فهو حر. وفي~$\lambd[x][x x]$ تكون واقعتا~$x$ مقيدتين
  بـ$\lambd[x]$، لأنهما حرتان في~$xx$. وفي~$((\lambd[x][xx])x)$ تكون
  الواقعة الأخيرة لـ$x$ حرة لأنها ليست في مجال~$\lambd[x]$. وفي
  $\lambd[x][(\lambd[x][x])x]$ يكون مجال~$\lambd[x]$ الأولى هو
  $(\lambd[x][x])x$، ومجال~$\lambd[x]$ الثانية هو الواقعة قبل الأخيرة
  لـ$x$. وفي~$(\lambd[x][x])x$ تكون الواقعة الأخيرة لـ$x$ حرة، والواقعة
  قبل الأخيرة مقيدة. ولذلك تكون الواقعة قبل الأخيرة لـ$x$ في
  $\lambd[x][(\lambd[x][x])x]$ مقيدة بـ$\lambd[x]$ الثانية، والواقعة
  الأخيرة مقيدة بـ$\lambd[x]$ الأولى.
\end{ex}

يمكننا، لحد~$P$، فحص جميع وقائع المتغيرات فيه والحصول على مجموعة
المتغيرات الحرة. ويرمز إلى هذه المجموعة بـ$\FV{P}$، ولها التعريف الطبيعي
الآتي:

\begin{defn}[المتغيرات الحرة لحد] \ollabel{def:fv}
  تعرّف مجموعة \emph{المتغيرات الحرة} لحد استقرائيًا كما يأتي:
  \begin{enumerate}
    \item $\FV{x} = \{x\}$ \ollabel{def:fv1}
    \item $\FV{\lambd[x][N]} = \FV{N} \setminus \{x\}$    \ollabel{def:fv2}
    \item $\FV{PQ} = \FV{P} \cup \FV{Q}$ \ollabel{def:fv3}
  \end{enumerate}
\end{defn}

\begin{prob}
  \begin{enumerate}
  \item عيّن مجالات~$\lambd[g]$ ونسختي~$\lambd[x]$ في الحد:
    $\lambd[g][(\lambd[x][g (x x)]) \lambd[x][g (x x)]]$.
  \item هل كل وقائع المتغيرات في
    $\lambd[g][(\lambd[x][g (x x)]) \lambd[x][g (x x)]]$ مقيدة؟ وبأي
    تجريد تقيد كل منها؟
    \item أعط~$\FV{\lambd[x][(\lambd[y][(\lambd[z][x y]) z]) y]}$.
  \end{enumerate}
\end{prob}

\begin{explain}
يشبه المتغير الحر إحالة إلى العالم الخارجي (أي \emph{البيئة})، ويمكن النظر
إلى الحد الذي يحتوي متغيرات حرة بوصفه حدًا محددًا جزئيًا، لأن سلوكه يعتمد
على كيفية إعداد البيئة. فمثلًا، في الحد~$\lambd[x][f x]$ الذي يقبل
وسيطًا~$x$ ويعيد ناتج تطبيق~$f$ على ذلك الوسيط، يكون المتغير~$f$ حرًا. وتعتمد قيمة الحد
على البيئة التي يوجد فيها، وبخاصة على قيمة~$f$ في تلك البيئة.

وإذا طبقنا التجريد على هذا الحد حصلنا على~$\lambd[f][\lambd[x][f
    x]]$. ولا يعود هذا الحد معتمدًا على متغير البيئة~$f$، لأنه يدل الآن
على دالة تقبل وسيطين وتعيد نتيجة تطبيق الأول على الثاني. ولن يكون لتغيير
$f$ في البيئة أي أثر في سلوك هذا الحد، لأن الحد لا يستعمل إلا ما يمرر
إليه وسيطًا، لا قيمة~$f$ في البيئة.
\end{explain}

\begin{defn}[الحد المغلق، المؤلف]
  يسمى الحد الذي لا متغيرات حرة فيه \emph{حدًا مغلقًا} أو
  \emph{مؤلفًا}.
\end{defn}

\begin{lem}\ollabel{lem:fv}
  \begin{enumerate}
    \item \ollabel{lem:fv-abs} إذا كان~$y \neq x$، فإن~$y \in
      \FV{\lambd[x][N]}$ إذا وفقط إذا كان~$y \in \FV{N}$.
    \item \ollabel{lem:fv-app} يكون~$y \in \FV{PQ}$ إذا وفقط إذا كان
      $y \in \FV{P}$ أو~$y \in \FV{Q}$.
    \end{enumerate}
\end{lem}

\begin{proof}
  تمرين.
\end{proof}

\begin{prob}
  أثبت~\olref[lam][syn][fv]{lem:fv}.
\end{prob}

\end{document}
